\begin{tikzpicture}[scale=0.62, every node/.style={font=\small}]
  % ============================================================
  % kji-space chunking, drawn on a single ji plane (k fixed).
  %
  % Grid convention: cell (i,j) occupies x in [i,i+1], y in [j,j+1].
  % Interior cells: i in 1..6, j in 1..4  (ni=6, nj=4).
  % Ghost ring of width nghost=1: i in {0,7}, j in {0,5}.
  % The full memory grid is therefore 8 (i:0..7) by 6 (j:0..5).
  %
  % Logical (flat/coords) order walks i fastest, then j: bottom row
  % left-to-right, then the next row up, ... , touching only interior
  % cells. Memory order is the same row-major walk but over the full
  % allocated grid, so a chunk that crosses a row boundary sweeps the
  % ghost cells that lie between logical cells in the contiguous span.
  % ============================================================

  \colorlet{cA}{blue!22}
  \colorlet{cB}{orange!38}
  \colorlet{cC}{green!28}
  \colorlet{cD}{purple!22}
  \colorlet{ghost}{gray!18}

  % Helper: draw the ghost ring (as light-filled border) for a panel
  % whose scope origin is the current (0,0). Interior box is (1,1)-(7,5).

  % ------------------------------------------------------------
  % Panel (a): i_pencil, logical -- one interior row per chunk.
  % ------------------------------------------------------------
  \begin{scope}[yshift=0cm]
    \node[font=\bfseries, anchor=west] at (0,6.4)
      {(a) \texttt{i\_pencil}, logical: one interior row = one chunk};

    % ghost ring
    \fill[ghost] (0,0) rectangle (8,6);
    \fill[white] (1,1) rectangle (7,5);

    % four row-chunks over the interior
    \fill[cA] (1,1) rectangle (7,2);
    \fill[cB] (1,2) rectangle (7,3);
    \fill[cC] (1,3) rectangle (7,4);
    \fill[cD] (1,4) rectangle (7,5);

    \draw[step=1, gray!55, very thin] (0,0) grid (8,6);
    \draw[very thick, gray!70] (1,1) rectangle (7,5); % interior boundary

    \node[anchor=west] at (8.3,0.5) {$i$};
    \node[anchor=south] at (0,6.0) {$j$};

    \node[cA!30!black] at (4,1.5) {chunk 0};
    \node[cB!55!black] at (4,2.5) {chunk 1};
    \node[cC!45!black] at (4,3.5) {chunk 2};
    \node[cD!55!black] at (4,4.5) {chunk 3};

    \node[gray!55!black, font=\scriptsize, anchor=west] at (8.3,5.3)
      {ghost cells};
    \node[gray!55!black, font=\scriptsize, anchor=west] at (8.3,4.7)
      {(never touched)};
  \end{scope}

  % ------------------------------------------------------------
  % Panel (b): ninner=8, logical -- chunks straddle row boundaries
  % but touch only interior cells.
  % ------------------------------------------------------------
  \begin{scope}[yshift=-8cm]
    \node[font=\bfseries, anchor=west] at (0,6.4)
      {(b) $n_{\mathrm{inner}}=8$, logical: chunks straddle rows, interior only};

    \fill[ghost] (0,0) rectangle (8,6);
    \fill[white] (1,1) rectangle (7,5);

    % chunk 0: j=1 i=1..6 (6) + j=2 i=1..2 (2)
    \fill[cA] (1,1) rectangle (7,2);
    \fill[cA] (1,2) rectangle (3,3);
    % chunk 1: j=2 i=3..6 (4) + j=3 i=1..4 (4)
    \fill[cB] (3,2) rectangle (7,3);
    \fill[cB] (1,3) rectangle (5,4);
    % chunk 2: j=3 i=5..6 (2) + j=4 i=1..6 (6)
    \fill[cC] (5,3) rectangle (7,4);
    \fill[cC] (1,4) rectangle (7,5);

    \draw[step=1, gray!55, very thin] (0,0) grid (8,6);
    \draw[very thick, gray!70] (1,1) rectangle (7,5);

    \node[anchor=west] at (8.3,0.5) {$i$};
    \node[anchor=south] at (0,6.0) {$j$};

    \node[cA!30!black] at (2,1.5) {chunk 0};
    \node[cB!55!black] at (5,2.5) {chunk 1};
    \node[cC!45!black] at (4,4.5) {chunk 2};

    \node[align=left, anchor=west, font=\scriptsize] at (8.3,3.4)
      {Each chunk is 8 interior\\
       cells in flat order\\
       ($i$ fastest, then $j$).\\
       Chunk 0 spills from\\
       row $j{=}1$ into row $j{=}2$.};
  \end{scope}

  % ------------------------------------------------------------
  % Panel (c): memory -- same chunk 0 as (b), but the contiguous
  % memory span also sweeps the ghost cells between the rows.
  % ------------------------------------------------------------
  \begin{scope}[yshift=-16cm]
    \node[font=\bfseries, anchor=west] at (0,6.4)
      {(c) memory: chunk 0's contiguous span also sweeps inter-row ghosts};

    \fill[ghost] (0,0) rectangle (8,6);
    \fill[white] (1,1) rectangle (7,5);

    % chunk 0 logical cells (solid)
    \fill[cA] (1,1) rectangle (7,2);   % j=1 i=1..6
    \fill[cA] (1,2) rectangle (3,3);   % j=2 i=1..2

    % swept ghost cells inside the contiguous memory span:
    %   trailing ghost of row j=1 (i=7) and leading ghost of row j=2 (i=0)
    \fill[red!22] (7,1) rectangle (8,2);   % j=1, i=7 (ghost)
    \fill[red!22] (0,2) rectangle (1,3);   % j=2, i=0 (ghost)

    \draw[step=1, gray!55, very thin] (0,0) grid (8,6);
    \draw[very thick, gray!70] (1,1) rectangle (7,5);

    % emphasize the two swept ghosts
    \draw[red!75!black, very thick] (7,1) rectangle (8,2);
    \draw[red!75!black, very thick] (0,2) rectangle (1,3);
    \node[red!75!black, font=\bfseries] at (7.5,1.5) {g};
    \node[red!75!black, font=\bfseries] at (0.5,2.5) {g};

    \draw[-{Latex}, red!75!black, thick]
      (6.5,1.5) .. controls (7.6,1.5) and (7.6,2.5) .. (0.5,2.5);

    \node[anchor=west] at (8.3,0.5) {$i$};
    \node[anchor=south] at (0,6.0) {$j$};

    \node[cA!30!black] at (3,1.5) {chunk 0 (logical cells)};

    \node[align=left, anchor=west, font=\scriptsize] at (8.3,3.4)
      {The memory span runs\\
       from $(j{=}1,i{=}1)$ to\\
       $(j{=}2,i{=}2)$ as one\\
       contiguous run, so the\\
       two ghost cells (\textcolor{red!75!black}{g})\\
       between the rows are\\
       swept as well. Logical\\
       cells are still each\\
       touched exactly once.};
  \end{scope}

\end{tikzpicture}
